% Issue #87: 添加书耳功能
% https://github.com/open-guji/luatex-cn/issues/87
%
% 书耳：刻在边栏外上角的小框，记篇名。\banxinSetup{书耳=true} 开启，
% 文字默认取当前 \chapter 的篇名（换篇后跟着变）；位置默认在左侧，
% 筒子页拆分后只落在左半叶那一页上，奇偶页自然区分。
%
% 本文件：两篇各占一个筒子页（四页），第 2、4 页（左半叶）的书耳分别是「五帝本紀」「夏本紀」；
% 篇名写在正文里的 \chapter，换篇用 \clearpage
\documentclass{ltc-guji}
\setmainfont{TW-Kai}
\banxinSetup{书耳=true}
\title{史記}

\begin{document}
\begin{正文}
\chapter{五帝本紀}
黃帝者，少典之子，姓公孫，名曰軒轅。生而神靈，弱而能言，幼而徇齊，長而敦敏，成而聰明。軒轅之時，神農氏世衰。諸侯相侵伐，暴虐百姓，而神農氏弗能征。於是軒轅乃習用干戈，以征不享，諸侯咸來賓從。而蚩尤最為暴，莫能伐。炎帝欲侵陵諸侯，諸侯咸歸軒轅。軒轅乃修德振兵，治五氣，蓺五種，撫萬民，度四方，教熊羆貔貅貙虎，以與炎帝戰於阪泉之野。三戰，然後得其志。
\clearpage
\chapter{夏本紀}
黃帝者，少典之子，姓公孫，名曰軒轅。生而神靈，弱而能言，幼而徇齊，長而敦敏，成而聰明。軒轅之時，神農氏世衰。諸侯相侵伐，暴虐百姓，而神農氏弗能征。於是軒轅乃習用干戈，以征不享，諸侯咸來賓從。而蚩尤最為暴，莫能伐。炎帝欲侵陵諸侯，諸侯咸歸軒轅。軒轅乃修德振兵，治五氣，蓺五種，撫萬民，度四方，教熊羆貔貅貙虎，以與炎帝戰於阪泉之野。三戰，然後得其志。
\end{正文}
\end{document}
